% Source: https://artofproblemsolving.com/community/c3h2210438_prove_the_obvious_problems
% Stable identity: aops-topic-2210438-q01
\textbf{(a)} If $a=b$, then $a^6$ and $b^6$ are the same number raised to the same power: $a^6=a\cdot a\cdot a\cdot a\cdot a\cdot a=b\cdot b\cdot b\cdot b\cdot b\cdot b=b^6$, replacing each factor $a$ by the equal number $b$.

\textbf{(b)} $a=1$, $b=-1$: $1^6=1=(-1)^6$, but $1\ne-1$.

\textbf{(c)} Factor:
\[a^6-b^6=(a^3-b^3)(a^3+b^3)=(a-b)(a+b)(a^2+ab+b^2)(a^2-ab+b^2).\]
Now $a^2+ab+b^2=\left(a+\dfrac {b}{2}\right)^2+\dfrac{3}{4}b^2$ and $a^2-ab+b^2=\left(a-\dfrac {b}{2}\right)^2+\dfrac{3}{4}b^2$ are positive unless $a=b=0$. So if $a^6=b^6$ and not $a=b=0$, then $(a-b)(a+b)=0$, i.e.\ $a=b$ or $a=-b$; and if $a=b=0$ this holds too. Conversely, $a=\pm b$ gives $a^6=(\pm b)^6=b^6$. $\blacksquare$
